\chapter{北京地铁4号线·大兴线}
\begin{epigraph}
    北京地铁4号线与大兴线贯通运营，大致呈南北走向，串联海淀、西城、丰台、大兴4个行政区，北起安河桥北站，南至天宫院站，全长50.7公里，设有35座车站，换乘站10座。
\end{epigraph}

\section{安河桥北}
安河桥北站位于海淀区中部，京密引水渠东侧，东临龙背村路。

初设地铁站时因紧邻龙背村而称“龙背村站”，后因南邻安河桥改称“安河桥北站”。龙背村临近安河桥，因肖家河（今清河西段）裁弯取直以前，该村邻近的河道向北呈弯曲形状，远远望去，形似龙背，故此得名。

作为古桥，安河桥已有200多年的历史。该桥的修建与清代圆明园护军营正红旗的设立有关。清雍正二年（1724年），为了保卫皇家园林圆明园的安全，清廷特设守卫圆明园及保卫皇帝由紫禁城赴圆明园沿途安全的“圆明园八旗内务府三旗护军营”，简称“圆明园护军营”，其正红旗在萧（肖）家河北侧，河南为丰益仓，由正红旗营兵看守。为通行方便，在河上修建了一座桥梁。最初是一座木桥，但因过往的人员车辆较多，没过几年木桥就被压坏，随后在原址上重修一座石桥，另在桥的附近设置水闸，控制上游水流，并取“安澜平和”之意，称“安和桥”，在桥拱正中央置石额，上刻“安和桥”三个大字。因石桥呈拱形，俗称“罗锅桥”，又因临近丰益仓而被称为“丰益桥”。

此外，“安和桥”还称“安河桥”，当时的萧家河是西山、玉泉山一带河湖的泄水河，此处设有水闸，其功能是在正常水情下向瓮山泊（今颐和园昆明湖）输水；在山洪暴发时则向清河排洪，所以“安河”乃“安定河水，防灾慰民”之意。

作为古村落，安河桥村大约在清光绪年间形成，因紧邻安河桥得名。民国时期已成较大的聚落，有安河桥西街一巷、二巷，安河桥北胡同、后胡同，礼拜寺胡同等。20世纪六七十年代为海淀人民公社青龙桥大队所辖村，而今已消失，地名尚存。

\section{北宫门}
北宫门站位于海淀区中部偏南，颐和园北宫门东侧。

北宫门是“三山五园”之一的颐和园北门，始建于乾隆年间，原名“北楼门”，光绪年间重建后改称“北宫门”。

北宫门坐南朝北，为颐和园的主要宫门之一，是一座面阔七间的两层门楼，歇山式，周围有廊。前檐两层明间悬匾额“凤策仰辉”，后檐两层明间悬匾额“兰馨菊香”，其门联为“雉扇开时娲簧喜奏齐天乐，凤韶谱处舜琯偕陈益地图”。

据传，北宫门在乾隆年间曾是清漪园（今颐和园）的正门，因清漪园南面为昆明湖，并没有大门，视北为上，故北侧的宫门为正门。乾隆至咸丰年间，帝后及大臣们到玉泉山静明园、香山静宜园，多出此门，经青龙桥镇一直向西，过功德寺至静明园等处。

1860年清漪园被“英法联军”焚烧，同治年间慈禧太后“垂帘听政”时，挪用建海军的款项重建，并更名为“颐和园”。同时将东宫门改为颐和园的正门。从此北宫门便成为颐和园的北门，其名气也远不如东宫门知名了。

进入北宫门是颐和园万寿山的后山，早年间有许多建筑。其中苏州街是仿江南水镇而建的，俗称“买卖街”。清漪园时期岸上有各式店铺，如玉器古玩店、绸缎店、点心铺、茶楼、金银首饰楼等。店铺中的伙计都是太监、宫女装扮。皇帝游幸时开始“营业”。后湖岸边的数十处店铺，1860年被列强焚毁。现在的景观为1986年重修，是颐和园后山的主要景区。

\section{西苑}
西苑站位于海淀区中部偏西，颐和园路与中直路交会处西北侧，可与16号线换乘。

《汉书》曰：“苑，谓马牧也。”远在明代，此地水草丰美、林木茂盛，适养牲畜、禽兽，故而称“苑”，因地处京城之西而得名“西苑”。关于这一地名的具体所指，有几种说法。

一是指颐和园东宫门外的街区及村落，曾有西苑一道街、西苑二道街和西苑三道街，其前街的操场为西苑之营盘，1915年至1928年设有西苑兵营。后因此地北靠圆明园，西邻颐和园，逐渐形成商业区，并持续到本世纪初。西苑地区的西苑营市街，曾是京西一处商业街区，繁荣一时。

二是泛称海淀一带的皇家园林。清代从康熙年间修建畅春园起，就有园居的习惯，此后又在西郊营造了清漪园（颐和园）、圆明园等多处皇家行宫御苑，既有宫廷建筑的雍容华贵，又有江南水乡园林的委婉多姿，皆因位于北京城的西北部而泛称“西苑”。

三是清康熙年间将畅春园称为“西苑”，位于今天的北京大学西侧，是康熙皇帝在“李园”旧址上建造的规模宏大的第一座皇家园林。“李园”本是明朝神宗的外祖父李伟修建的“清华园”，因建筑规模宏大，曾被称为“京师第一名园”，后废弃。康熙二十三年（1684年），康熙皇帝南巡归来后在此大兴土木。康熙二十九年（1690年），新修建的园林被命名为“畅春园”。

咸丰十年（1860年），“英法联军”攻入北京焚烧圆明园时，将其一并烧毁。此后畅春园旧址失于保护，园内残存建筑在同治年间被拆，用于圆明园复建工程。光绪二十六年（1900年）畅春园又遭“八国联军”洗劫，至民国时期，其遗址已成荒野，仅有恩佑寺和恩慕寺两座山门残存，即今天北京大学西门外南侧的两处古建筑遗迹。

\section{圆明园}
圆明园站位于海淀区中部偏西，清华西路中段，圆明园遗址公园南门西侧。

圆明园由圆明园、长春园、绮春园（后改名为万春园）组成，呈倒“品”字形。三园紧相毗连，统称“圆明园”，始建于康熙四十八年（1709年），最初是康熙皇帝赐给四子雍亲王胤禛的赐园。雍正即位后将其拓展，御以“避喧听政”。乾隆皇帝在位期间除对圆明园进行局部增建、改建之外，还在东邻新建了长春园，在东南邻并入了万春园，由此圆明三园的格局基本形成。它是自康熙末叶起，历经150余年建造的一座大型皇家宫苑。雍正、乾隆、嘉庆、道光、咸丰五代皇帝每年大部分时间居住在圆明园，在此举行朝会，处理政事。

“圆明园”为康熙命名，并御书金字匾牌悬挂于圆明园殿门楣上方。对“圆明”二字，雍正皇帝的解释是：“圆而入神，君子之时中也；明而普照，达人之睿智也。”另说康熙题名“圆明园”，是取意于雍正信奉道教，其法号为“圆明”。

1860年“英法联军”焚毁了圆明园，同治帝时欲修复，后因财政困难被迫停止。圆明园1900年再次遭到“八国联军”的破坏，此后又遭到军阀、地痞等的巧取豪夺，残存的建筑被拆毁，石料、木料被掠夺，逐渐成为一片废墟。随后，附近的村民及来自河北、河南的流民在园内任意开垦，平山填湖，种植水稻和蔬菜，就近建房居住，短短几十年就形成了大村落。20世纪80年代初属海淀人民公社二大队的自然村，因地处圆明园内而称“圆明园村”。

1983年《北京市城市建设总体规划方案》明确规定了“整建圆明园……建成圆明园遗址公园”，随后成立了“圆明园遗址公园筹建委员会”，并对圆明园遗址部分景观恢复建设，逐渐将园内的村民、企业迁出。1988年6月29日，圆明园遗址公园正式对社会开放。

北京大学东门站位于海淀区东南部，中关村北大街与成府路交会处，北京大学东侧。

该站因地处成府路西口，初设置地铁站时称“成府路站”。这里最早是明代成王的墓地，以“阴间地府”称“陈府”，明朝有两位成亲王，一位是朱让栩，一位是朱弥，此陈府到底是哪位的墓地已无从考证。到了清朝，乾隆皇帝的第十一个儿子永瑆也被封为成亲王，所得宅园离陈府不远，故称“成亲王府”，只因这个王府是“阳宅府第”，永瑆于嘉庆元年（1796年），取消了“陈府”的称谓，改称“成府”。其具体位置在当时静春园东南，正蓝旗护军营房之西，南边与治贝子园搭界，这一带形成村落后，称“成府村”。1860年“英法联军”火烧圆明园时“成府”被焚毁。成府村有近30条胡同，成府街为其中之一，后向东延伸的道路称“成府路”。1982年至1985年清华园以西仍称“成府路”，以东因临近当时的东升人民公社驻地而称“东升路”。本世纪初该路拓宽、取直并向东延至京昌公路后，统称为“成府路”。

北京大学创建于清光绪二十四年（1898年），初名“京师大学堂”，校址在地安门内马神庙，1912年改名为“北京大学”，1916年迁校址至沙滩，后曾几次更名，1929年恢复“北京大学”校名。1952年全国高等学校院系调整，位于成府村西部燕园的“燕京大学”被撤销，其校舍由北京大学接收。最初并没有东门，只有一条出入校园的通道，后辟此门，门外紧邻成府村。

燕京大学创办于1919年，创始者为美国人司徒雷登先生。选址于海淀镇北部，以早年的淑春园遗址为中心，1926年基本落成。

\section{北京大学东门}

\section{中关村}
中关村站位于海淀区东南部，北四环西路与中关村北大街、中关村大街交会处南侧。

关于“中关村”的由来，追溯寻源，说法不一。

一说这里曾是永定河故道，有旱河流过，被称为“中湾”，形成村落后，称“中湾村”，以后河道逐步消失，便把“中湾村”谐音成了“中关村”。

二说清代有位太监在这里购置田地，打算作为养老和葬身之地。没过几年他就出宫住到这里，以后陆续来了不少年老的太监，于是这里被称为“太监营房”，后来有太监死了，就埋在附近。“中官”二字在过去有太监的意思，时间长了，这一带就成了“太监义地”，形成村落后，也称“中官坟”，后因忌讳“坟”字而讹称为“中官村”，又谐音为“中关村”。20世纪五六十年代在此开发建设时，曾挖出数十座太监坟墓。

三说清代雍正年间，在蓝旗营控制的辖区内设五关，其中蓝旗营汛防居中，所设之关被称为“中关”，以扼住西南通向东北的要道。此后中关一带形成村落，被称为“中关屯”，后改称为“中关村”。

四说光绪三十一年（1905年），慈禧太后七十大寿，京城搭了许多大戏台，连唱数日，以示祝寿。在海淀镇东北方向也一连搭了三座大戏楼，因高达两丈多，故称“小城关”。其北城关在圆明园南门的成府村外，南城关在今天的海淀黄庄附近，中城关在南北两个城关之间，也就是现在的中关村一带。后来人们就把搭建中城关的地方称为“城关村”，又因是中城关，也称“中关村”。

20世纪六七十年代，这一带属于东升人民公社辖区。80年代初，中关村附近科研单位的一批科技人员率先突破传统观念，从科研院所走出来，“下海”创办民营高新技术企业，将科研项目转化为科技成果。短短几年，一大批高新技术企业应运而生。这些企业大多沿中关村道路而经营，到1986年底，从白石桥到中关村之间的道路两侧已有近150家民营科技企业，形成了著名的“中关村电子一条街”，后发展为“中关村科技园区”。

“中关村”到底何时因何而确切得名，仍有待考证。2002年5月10日在中关村大街矗立起一块高8米的“中关村纪念碑”，给予中关村新的“定位”。

\section{海淀黄庄}
海淀黄庄站位于海淀区东南部，中关村大街与海淀南路、知春路交会处，可与10号线换乘。

“黄庄”本称“皇庄”，源于明洪武年间，发展于永乐初年，迅速增长于弘治年间。明朝时的土地分为官田和民田，官田最初叫“皇庄”“宫庄”，是皇室成员所经营的庄园，遍布京畿，占尽膏腴之地。嘉靖初年（1522年）明朝财政出现危机，世宗朱厚熜“改皇庄为官田”，从此京畿不复有皇庄。皇庄虽废，但与之性质相似的官庄一直保存到明末，而作为地名却流传到清代。1911年“辛亥革命”以后，取消封建帝制，以“皇”字相称的许多地名被改用“黄”字代替，“皇庄”也就改称了“黄庄”。

位于海淀镇东南部的黄庄，因临近西土城，也称“土城皇庄”。清代《日下旧闻考》称：“皇庄距南海淀二里许，盖沿明代俗称也。”清光绪年间出版的《京城简图》中标有“海甸皇庄”的字样，民国时的《京城街巷名称录》中已称“黄庄”，为区别于其他方位的“黄庄”，多称其为“海淀黄庄”。

20世纪50年代在海淀区划分行政区域时，海滨黄庄东部为东升人民公社，西部为海淀人民公社，以种植蔬菜为主。从80年代开始，这一带被迅速开发，如今已无农田。

\section{人民大学}
人民大学站位于海淀区东南部，北三环西路与中关村大街、中关村南大街交会处。

初设地铁站时因临近双榆树而称“双榆树站”，后因地处中国人民大学东南侧改称“人民大学站”。

据传，“双榆树”本称“桑榆墅”，其得名与清代权臣纳兰明珠有关。

纳兰明珠，叶赫那拉氏，满洲正黄旗人，康熙朝最重要的大臣之一，官居内阁13年，“掌仪天下之政”，故有“相国”之称。

康熙初年，明珠在此修建了一座别墅，占地百余亩，园中有庭舍、回廊、小楼、果木，尤以桑树、榆树为多，故取桑榆为别墅之名，称“桑榆墅”。

明珠当初建园于此，一则因为附近有明珠家的一处墓地；二则从京城到西郊海淀、畅春园等地，这里是必经之路，可用作中途休息之所；三则其长子纳兰性德喜欢辞赋，为他提供一个幽静之所，便于他读书。纳兰性德特别喜欢这里的景致，常在此与文士、好友谈论学问，吟诗唱赋。百余年后，此园的景致逐渐消失，成为村落，地名由“桑榆墅”谐音为“双榆树”。20世纪六七十年代这里建成居民区，称“双榆树小区”。

1950年中国人民大学由河北省正定县迁至双榆树村西侧，其前身是1937年在延安成立的陕北公学，以及后来的华北联合大学和华北大学。1950年10月，以华北大学为基础合并组建了中国人民大学，该校校址原是海淀镇南部的农田及乱坟岗子，零星地有几个村落，其西北部有一处坟圈子，据传是明末太监王承恩的墓地。该墓直到民初尚在，民国时此地形成聚落，时称“王公坟”。1952年兴建人民大学时村子被征用，王公坟被拆除，在其西侧建成万泉庄小区，村民多搬迁至此，今称“王公坟小区”。

\section{魏公村}
魏公村站位于海淀区东南部，魏公村路、学院南路与中关村南大街交会处北侧。

魏公村的来历和古代畏兀儿人（维吾尔族）有关，目前通常的说法是此地因元朝初期为畏兀儿人聚居地而得名“畏吾村”，明代称“苇孤村”，清代称“魏吴村”“卫伍村”，后谐音为“魏公村”。

明代成化年间，太监墓志铭称这个村为“谓务村”。正德年间，太监墓志铭称“委兀村”，张爵《京师五城坊巷胡同集》称“畏吾村”，万历年间成书的《宛署杂记》称“苇防护林地”。直到清道光十四年（1834年），这个村子仍叫“畏吾村”。清人乔松年《萝蘼亭札记》载：“畏兀村，盖京西直门外村名。本西域畏兀部落，元太祖时来归，聚处于此，以称村焉。”清人查礼《畏吾村考》中写道：“京师西直门外，有村名畏吾。”1915年绘制的《实测京师四郊图》标有“魏公村”。“文革”期间改为“为公村”，1982年又恢复“魏公村”之名。

《魏公村研究》称，魏公村在元朝时是畏兀儿人仕官的宅墅及其家族墓地之所在，如受封宅邸并敕封“魏国公”的阿里海涯及其子廉希宪家族都葬于此。元帝还追封畏兀儿重臣扎马拉丁、阿里罕、亦不拉金祖孙三人为“魏国公”。因此推断，魏公村并不像有些学者所考证的是讹称、误读，更不是如“魏吴”之舛写，而是因为廉希宪等畏兀儿人被追封为“魏国公”，清代此地形成聚落称“魏公村”。

魏公村可考的历史，从元朝开始。忽必烈的重臣、他的连襟、畏兀儿名臣、前燕京行省札鲁忽赤，蒙速思死后葬在这里。后来，另有两位畏兀儿名臣廉希宪、阿里海涯也葬在高梁河畔。由于这些畏兀儿大族坟茔的存在，这里形成了由畏兀儿守墓坟户和农户组成的村落，称之为“畏兀村”，后谐音为“魏公村”。

\section{国家图书馆}
国家图书馆站位于海淀区东南部，国家图书馆东侧，南邻白石桥，可与9号线换乘。

因地处白石桥北侧，初设地铁站时称“白石桥站”。白石桥始建于元至元二十九年（1292年），横跨高梁河（长河），因使用的石料为白色，时称“小白石桥”，明代重建后称“白石桥”，清代改建为单孔、米黄色花岗岩砌筑的梯形石桥，依然称“白石桥”。桥长10米，宽6米，高3米，两侧立有石刻护栏，两端戗以抱鼓石，是北京地区保存较好的古桥之一。

清代时，此处也称“郑亲王坟”，又称“王爷坟”，坟主人是清顺治皇帝之叔济尔哈朗，被封为郑亲王，后代世袭，最后一位郑亲王叫端华，即清光绪年间军机大臣肃顺的哥哥，后来被慈禧太后赐死。在郑王坟旁边有看坟人居住的地方，叫“胡家楼”，后来展宽西颐路（今中关村南大街），“胡家楼”地名遂取消，唯有“白石桥”这个古老的地名被保留下来。

为适应现代交通的需要，1982年在古桥西侧又建了一座钢筋混凝土梯形跨河公路桥，两桥“合二为一”。1998年扩建白石桥路时，拆除了旧石桥的护栏，在扩展后的两侧修建了白色的新栏杆。

“国家图书馆”原称“北京图书馆”。1909年9月，清政府批准军机大臣、学部尚书张之洞的奏请，筹建“京师图书馆”。1916年教育部饬京师图书馆，凡在内务部立案的出版图书均交京师图书馆庋藏。1949年改名为“北京图书馆”，原址在北海附近的文津街，1987年在白石桥西北侧建成新馆，被确定为中国唯一的国家级图书馆，1998年12月更名为“国家图书馆”。

\section{动物园}
动物园站位于西城区西北部，西直门外大街中段，北京动物园南侧。

北京动物园的历史可追溯到清光绪三十二年（1906年），当时被称为“万牲园”，其前身为清农工商部农事试验场，集农作物种植试验与动物养殖于一体，兼有公园性质，是在乐善园、继园和广善寺、惠安寺的旧址上兴建的。

光绪三十四年（1908年）6月，农事试验场全部竣工，开放接待游人。最初展览的动物是两江总督端方自德国购回的部分禽兽及全国各地抚督送献清朝政府的动物，约百余种。

辛亥革命后，农事试验场几易其名，从“中央农事试验场”到“国立北平天然博物院”再到“实业总署园艺试验场”直至“北平市园艺试验场”。由于连年战乱，民生凋敝，大部分动物夭折，所剩无几。1949年9月1日定名为“西郊公园”，将园内的围墙、牡丹亭、豳风堂、动物园兽舍等加以修缮，并整修了鸟笼、鹿棚、猴山和甬路，增添一些小型动物。

随着动物种类和数量的增加，动物园的基本设施建设也得到迅速发展。1955年4月1日，西郊公园正式改名为“北京动物园”，此后兴建了象房、狮虎山、猕猴馆、猩猩馆、海兽馆、两栖爬行动物馆等场馆，其中狮虎山、猩猩馆、两栖爬行馆等场馆，使用至今，并成为北京动物园的标志建筑。

而今北京动物园已饲养展览动物500余种，5000多只；海洋鱼类及海洋生物500余种，10000多尾，成为中国开放最早、饲养展出动物种类最多的动物园，也是中国最大的城市动物园和世界知名的动物园。

\section{西直门}

\section{新街口}
新街口站位于西城区中部偏北，西直门内大街与赵登禹路、东新开胡同交会处。

此地元代曾有漕运水道与积水潭相连。明代改建元大都时，积水潭码头被废弃，将西直门北侧至今新街口北边的水域填成陆地，有的建为居民区，有的修成道路。相对于旧有街道，便有“新开”之意，称“新开道街”，后衍化成“新街口”，即新开街道的路口。明嘉靖年间《京师五城坊巷胡同集》中便出现了“新开路”的地名，而万历年间的《宛署杂记》已有“新街口”的记载。清代昭梿《啸亭续录》有“董太保赐地新街口，谦郡王府在羊肉大街”之记载。明代后期，新街口地区的工商业已发展起来，从今天的新街口东街西口往南至护国寺大街西口，店铺一家挨一家。探其原因，一是新街口靠近西直门，为进出西直门的必经之地；二是护国寺庙市对新街口地区工商业繁荣影响极大。清末和民国年间这里成为京城的商业中心。

新街口地处丁字路口，早年间往西可看西直门城楼，往北能望到北城墙，往南可到西四牌楼。由路口往西到崇元观，也就是现在的赵登禹路北口，称“新街口西大街”，往北到城根称“新街口北大街”，往南到护国寺称“新街口南大街”。

20世纪50年代进行城市改造时，为方便车辆出入，在新街口北面的城墙上打开一个缺口，使北京城西北角又多了一个通道，称“新街口豁口”。出豁口往北，是一大片水域，称“太平湖”，原本是和积水潭连成一片的。随着豁口的打开，新街口一带更加繁荣，是京城西北部重要的商业街区，其南北大街两侧商店众多，以日用百货、服装、餐饮等综合商业为主。

\section{平安里}
平安里站位于西城区中部，平安里西大街、地安门西大街与新街口南大街、西四北大街交会处，可与6号线换乘。

“平安里”这一地名只有80多年的历史。相传明朝时此地为太平仓，平安里原本是明代太平仓后面的狭窄无名夹道，处于皇城的西北角之西。太平仓建于明正德五年（1510年），前身为永昌寺。东起今护仓胡同，西抵西四北大街、新街口南大街，北始群力胡同，南达太平仓胡同。清代废弃，改为承泽亲王府。承泽亲王是顺治第五子硕塞，其子博果铎袭封，改号庄，无嗣，康熙皇帝以十六子允禄为后，其府遂称庄王府。1900年“庚子之变”中毁于大火。据说曾有人在此挖出黄金。为了获得财宝，民国军阀李纯与其弟李馨，仅以20万银元就从穷困潦倒的末代庄亲王手中买走了这座王府。李纯系直隶（今天津市）人，曾任江西都督、江苏督军。在任期间于津京两地广置房产，是当时天津最大的房产主之一。但他在拆除王府的旧建筑时并没有挖出什么财宝，遂将庄王府的地上建筑尽数拆除，把砖瓦、木料运到天津建了李氏祠堂（今南开文化馆）。此后在此重新建房，其样式为中西合璧式房屋，门前无名夹道被拓宽，并取了一个吉祥的名字——平安里。1938年出版的《最新北平全图》中已有了“平安里”的街名，此后其附近的街区也被泛称为“平安里”。

北京历史上曾有五个叫“平安里”的地方，分别在东城、西城、宣武、丰台和昌平，但只有位于西城的平安里最为著名，并保留至今。

\section{西四}
西四站位于西城区中部偏北，阜成门内大街、西四东大街与西四北大街、西四南大街交会处。

“西四”是“西四牌楼”的简称，处于东、西、南、北相通的十字路口，早年间每个路口各建有一座四柱三楼式木结构牌楼，檐下有如意斗拱，俗称“四牌楼”或“西四”。

据《日下旧闻考》载：“宣武门北有单排楼曰瞻云，又北二里有四牌楼，东曰行义，西曰履仁，南北曰大市街。”故西四牌楼南北路口的牌楼上书“大市街”，东路口牌楼上书“行义”，西路口牌楼上书“履仁”。它与皇城以东的东四牌楼相对称，是北京城两个繁华的商业街。

西四牌楼以北，曾叫“西大市街”，是南北走向的街道。西四牌楼以南，从明代至清光绪年间也称“西大市街”，清宣统时称“丁字街”。因东为西安门大街，往西没有街巷，形成“丁”字。西四牌楼东面的街道很短，往东约三百米就是皇城根。西四牌楼西面一直到阜成门，从明代至清乾隆年间，称“阜成门大街”。约在乾隆年后，从沟沿往东，直至西四牌楼称“羊市大街”；沟沿以西，仍称“阜成门大街”。

西四十字路口，是明代处决犯人的刑场。英宗复辟后，天顺元年（1457年）保卫北京的于谦也在这里被杀害；武宗时期的大太监刘瑾在这里被处死；崇祯二年（1629年），后金军队围困北京，袁崇焕带兵增援大败金兵，崇祯皇帝中了皇太极的反间计，袁崇焕在这里含冤被杀。清朝则将该刑场改到宣武门外菜市口。

1954年12月底到1955年1月中在展宽马路时，将西四牌楼拆除，但“西四”这一地名沿用至今，成为京城繁华的商业街区。

\section{灵境胡同}
灵境胡同站位于辟才胡同、灵境胡同与西单北大街交会处北侧。

明朝时，灵境胡同分东西两部分，东段因坐落有灵济宫，被称为灵济宫；西部南侧有宣城伯府，故称宣城伯后墙街。清朝时，以西黄城根南街为界，东段因原“灵济宫”逐渐谐音为“灵清宫”“林清宫”（传说与嘉庆年间攻打皇宫的义民林清有关），因此被称为“林清胡同”，西段则称为“细米胡同”。1911年后，东段改称“黄城根”，西段则称为“灵境胡同”。1949年后两段并称为“灵境胡同”。

灵济宫全称“洪恩灵济宫”，是一座道教宫院，供奉的是洪恩灵济真君，俗称“二徐真君”，他们分别是五代吴国大臣徐知谔、徐知证，以“宽仁爱物，同心好善，精勤至道”而著称。

据《帝京景物略》载：明永乐十五年（1417年），朱棣患病，入睡时，梦见两位道士前来敬献灵丹妙药，结果没过三天，他的病就痊愈了。为了感谢仙道的救命之恩，朱棣下令在城中为其建宫祭祀，封其为玉阙真人和金阙真人，赐名“洪恩灵济宫”。因是皇家敕建的庙宇，其规模相当宏大。每逢初一、十五、立冬、夏至等节令，皇家总要派大臣们前来烧香祷告，祭祀真人。

崇祯十五年（1642年），有位大臣向皇帝写了一个奏章，说灵济宫供奉的两位真人是叛臣之子，不宜受朝臣敬拜，请示用帷帐将塑像遮盖起来，停止祭祀活动。崇祯皇帝信以为真，立即下令查封灵济宫，从此，香火旺盛了200多年的灵济宫衰落了。

20世纪80年代和90年代灵境胡同曾先后被改造，而今成为长664米，最宽处达32米的街道，已完全看不出胡同的模样了，俨然是一条现代化的大街，故有“北京最宽的胡同”之称，而“灵境胡同站”是北京众多地铁站名中唯一以“胡同”命名的。

\section{西单}
\section{宣武门}

\section{菜市口}
菜市口站位于西城区南部，广安门内大街、骡马市大街与宣武门外大街、菜市口大街交会处，可与7号线换乘。

早在一千多年前的辽代，这里是辽中都安东门外的郊野，金代为施仁门内的丁字街，明代是京城较大的蔬菜市场，沿街菜摊、菜店众多，故菜市最集中的街口易名“菜市街”，清代时改称“菜市口”，此名一直沿用至今。

菜市口名声大振的主要原因，是清政府将杀人的刑场从明朝时的西四牌楼（当时称西市）移至宣武门外的菜市口。据说当年的刑场就设于今天的菜市口大街北侧的十字路口附近。马芷庠所编《北平旅行指南》载：每逢秋后朝审，在京处决犯人众多之时，由东向西排列，刽子手执刀由东向西顺序斩决，由此“菜市口”成为“刑场”的代名词。问斩时监斩官的高座位常设于老字号“西鹤年堂”店门口。搭一席棚，下放一长方桌子，上摆朱墨、锡砚和锡制笔架，笔架上搁放几支新笔。监斩官照例先要在西鹤年堂坐一坐，稍事休息，再升座行刑。

清咸丰十一年（1861年），慈禧太后通过“辛酉政变”夺取大权，实行首次“垂帘听政”，而受咸丰皇帝遗诏的八位赞襄政务大臣中的肃顺，就是在此被杀头的。光绪二十四年（1898年）慈禧下令逮捕的“戊戌六君子”谭嗣同、林旭、杨锐、杨深秀、刘光第和康广仁也被杀害于此。

1911年随着清王朝的灭亡，刑场被转移至郊外，此后这一带逐渐成为繁华的商业街。后经改扩建，如今东为骡马市大街，南为菜市口大街，西为广安门内大街，北为宣武门外大街，成为京城西南部的交通枢纽。

\section{陶然亭}
陶然亭站位于西城区东南部，白纸坊东街、陶然亭路与菜市口大街交会处，陶然亭公园西北部。

据《陶然亭公园志》载：陶然亭为京师南城旧景，建于清康熙三十四年（1695年）。当时这一带是生产修建皇城所用砖瓦的窑厂，窑厂附近有一座元代古庙，称慈悲庵。工部郎中江藻是窑厂的监工，经常来此察看。他见这里有花有草有水塘，景色很好，便在慈悲庵的西面修建了三间宽敞的休息厅，并以自己的姓氏得名“江亭”。虽然得名为“亭”，其实就是供人歇息的房舍，并非园林景观中的亭子，故有“江亭无亭”之说。没过多久，江藻觉得这个名字太直白，有些俗气，于是翻阅了许多古代典籍，最后被唐代诗人白居易的诗《与梦得沽酒闲饮且约后期》最后一句“更待菊黄家酝熟，共君一醉一陶然”所吸引，觉得“陶然”二字别有意境，最富情致，于是将亭子改名为“陶然亭”。从此以后，他经常邀请一些文人雅士来此吟诗唱和，“陶然亭”也由此而出名。

慈悲庵是创建于元代的古刹，距今已有700余年历史。山门朝东，整个建筑布局严谨，瑰丽庄重。庙内西侧的三间敞轩就是人们常说的陶然亭。主要建筑有观音殿、准提殿、文昌阁、陶然亭等。自清代以来，这一带由于富有自然风光，一直是文人墨客聚会游览的地方，留下了很多诗文作品。

20世纪40年代，原有的部分建筑已经倒塌，只剩下遗迹，1949年后得以修复。从1985年开始，在此仿建了全国各地的十余座名亭，使陶然亭成为一座以亭文化为特色的公园。

\section{北京南站}
北京南站位于丰台区西南部，可与14号线换乘。

北京南站兴建的历史可以追溯到清末。光绪二十二年（1896年），北京至奉天（沈阳）的铁路北京至天津段通车，在永定门外马家堡村设站，称“马家堡站”，是北京最早的火车站。该站为英国人监造，有着典型的英式建筑风格。车站建成后，在其周围出现了很多新兴的店铺，例如茶棚、缸店、旅店、澡堂子、落子馆等，人来人往，热闹一时。1900年6月12日该站被“义和团”焚毁，1902年经英国人简单修复后恢复使用，因北距永定门不远，改称“永定门站”。因“八国联军”入侵北京而西逃一年多的慈禧太后回銮时，军机大臣便安排她从保定搭乘火车回京，在马家堡火车站举行了盛大的迎驾仪式。1906年此段铁路经永定门延伸到正阳门（前门），该站被撤销。

1957年，铁路部门决定在崇文门东侧修建新的北京火车站，为缓解施工期间北京地区火车运输压力，决定修建临时的永定门客运站，也就是在距离老永定门火车站以西1公里左右的地方修建新的永定门火车站。该站本打算作为临时客运站使用，设计寿命只有10年，结果超期使用了38年，1988年更名为“北京南站”。2006年5月进行改造，2008年8月开通运营。

北京南站所处的位置历史上为古村落彭庄，1949年以前有几十亩菜地，最初只有彭、贾两户，据传是清代乾隆年间从山东逃荒而来，在此安家，因彭为大姓户且居此年久，故称“彭庄”。当时附近除了菜地外，还有许多坟地和大苇塘，一条小路通往东北方向的永定门。1949年后，这一带的坟地和大苇塘经过平整建起了许多居民住房。1957年修建永定门火车站后，“彭庄”之名被逐渐淡化，但作为地名仍在使用，今属永外街道彭庄社区。

\section{马家堡}
马家堡站位于丰台区西北部，南三环西路与马家堡西路交会处南侧。

因地处京南交通要道，明末清初时今天的马家堡一带便出现了许多驿站、店铺，其中开店最早、铺面最大的是马家铺。

据传，开设马家铺的是回民马氏兄弟，他们为人和善，饭菜做得别有风味，更重要的是价格公道，生意做得很红火。短短几年，便有了些名气。日久天长，便把这里称为“马家铺”。十几年之后，他们有了一大笔积蓄，便在此修建了一座小楼，将小饭铺改成了既能吃喝又能住宿的客栈。清乾隆年间，在此开设店铺的人越来越多，逐渐形成较大的聚落，因马氏兄弟最早在此开设店铺，故称“马家铺”，后被谐音为“马家堡”。因“铺”与“堡”字义相近，均为驿站之义。

马家堡地处京城南边，清代时紧临南苑皇家园林围墙，马家堡在围墙西北部。南苑共有十三座角门，马家堡角门旧址在马家堡南街西口。南苑的角门都已无存，只留“马家堡角门”（今简化为“角门”）这一地名。

马家堡是北京城最早的火车总站。1896年，清政府修建卢汉（卢沟桥至汉口）铁路，1897年这条铁路由丰台接轨至马家堡，在距离永定门3公里的马家堡建造客货混运车站。车站气势恢宏，为当地地标式建筑。其实火车总站设在这里，还另有原因。当时火车还是个新鲜玩意儿，不但体积大噪声也大，慈禧害怕铁路进城破坏了北京城的“风水”，压着“龙脉”，原本铁路要修到正阳门（即前门），所以修到马家堡就不许再往前修了，于是马家堡成了北京最早铁路线的终点站。

作为北京的铁路总站，马家堡火车站虽然只存在了7年，但对于促进中国铁路的建设和发展起到了重要作用。

\section{角门西}
角门西站位于丰台区东部，嘉和路与马家堡西路交会处，可与10号线换乘。

该站因地处角门西侧而得名，角门实为马家堡角门，是清代皇家苑囿南海子（南苑）十三个角门之一，原址在马家堡南街一带，其迹已无存，只留下地名，而今成为京南地区的居民聚集地。

早在金代，南苑一带就是皇家禁苑，只是当时规模较小。元代为“下马飞放泊”，是帝王及贵族狩猎、游娱之地。明代成为北京城南一座风光绮丽的皇家苑囿。清乾隆年间又在南苑大兴土木，重建海子，将过去的土围墙改建成砖墙，用五尺厚三合土作根基垒四十八层灰色砖，墙基宽五尺，顶宽二尺，全长120多里，还将海子四门增为九门，除此之外还开辟了马家堡、栅子口、马道口等十三个角门。所谓“角门”，泛指官署、邸宅、园囿墙上小的旁门，均处于把角儿处。

马家堡角门位于永定门外西南六里处。实际上角门是南苑九门之外的便门，只有一个小门和一间值差（值班）的更房，是专为附近佃户（海户）、杂役进苑值差出入所设。

到了清末，列强闯入南苑，使其遭到严重破坏，并逐渐荒废，围墙、宫门及角门也随之消失，只留下一些与之相关的地名，角门是其中之一。民国时该地渐成农田并住有散居农户，随着村域逐渐扩大，被称为“角门村”。20世纪70年代后期，陆续在角门村及周边建成角门东里、角门路、角门北路等。

\section{新宫}
新宫站位于丰台区东南部，南苑西路与槐房西路交会处南侧。

“新宫”为“新衙门行宫”的俗称，是南苑四大行宫之一（其他三个是旧衙门行宫、南红门行宫和团河行宫，均已无存）。清代《日下旧闻考》卷七十五《国朝苑囿南苑二》记载：新衙门行宫位于南苑西北隅，“镇国寺门内五里许”。

据《南苑史话》载：新衙门曾为明代南海子内提督官署房，清朝初年改建成行宫，康熙年间名为“西宫”。新衙门行宫有宫门三楹，宫门前有元朝延祐元年（1314年）十月铸的一对铁狮子。左右垂花门里对面房十间，御题“迩延野绿”匾额。前殿三楹，后殿三楹。后殿屏扆间原有《乾隆大阅图》，后殿东有裕性轩五楹，是乾隆帝的书屋。裕性轩东小室的墙上，曾悬挂顺治帝御笔唐李白《清平调》诗，故苑户们将该小室俗称为“诗句房”。裕性轩西是澹思书屋，后有陶春室。另外还有古秀亭、春望楼等建筑。

1900年“八国联军”侵入南苑后，新衙门行宫遭到破坏，遂被废弃。民国时新衙门行宫建筑又遭到奉系军阀的破坏，所有建筑消失，附近逐渐发展成村落，称“新宫村”。20世纪80年代为南苑乡辖域，与东面的旧宫相对应。

新宫村为南苑乡南片行政村，由4个自然村组成，东临槐房村，西临花乡新发地，南临西红门镇，北临花乡草桥。2010年拆迁，建成新宫家园等小区。

\section{公益西桥}
公益西桥站位于丰台区东部，公益西街、枫竹苑北路与马家堡西路交会处。

公益西桥因临近公益庄而得名。公益庄地处明清皇家御囿南苑，清末列强闯入南苑后，明清两代苦心经营的南苑一蹶不振。光绪二十七年（1901年）七月《辛丑条约》签订后，为弥补国库空虚，于光绪二十八年（1902年）六月设立了“南苑督办垦务局”，将“南苑内闲旷地亩招佃垦荒”，实际上就是对外出租。南苑开放后，许多官僚、军阀、巨商、太监等纷纷抢占苑内的土地，建立私家庄园，龚记庄是其中之一。

据说，龚记庄是清朝宣统年间一位姓龚的四品官员的私家庄园。最初为一位王爷所有，那位王爷败落后转给这位龚官员，由此他雇用数十人在此耕种，并兴建了一座四合院，以为夏日避暑之所，俗称“龚记庄园”。随着庄园面积的扩大，所雇用人员不断增加，逐渐形成一个小的村落，称“龚家庄”或“龚记庄”。民国时该园逐渐废弃，而农田被当地村民继续耕种，于是“龚记庄”被谐音为“公益庄”，其旧址在今天公益桥南大红门火车站附近。

民国时期，公益庄属大兴县，实为大兴县与北平市三郊区接合处，有一条自北向南的大道，北至马家堡，南至新宫，道路两侧为大面积的农田。20世纪七八十年代这里是南苑人民公社的一个生产队。2000年修建的南四环路从村子北部经过，并分别在与马家堡路、马家堡东路和马家堡西路的交会处兴建了三座立交桥，于是将临近村子的立交桥称为“公益桥”，偏东的称“公益东桥”，偏西的称“公益西桥”。

\section{西红门}
西红门站位于大兴区北部，宏福西路、宏福路与欣宁街交会处南侧。

西红门一带辽金时形成村落，始称“西綦里”。元代在该村东部建“下马飞放泊”，供皇家狩猎之用。明代湖泊面积逐渐扩大，更名为“南海子”，且大兴土木，建成一处规模宏大的皇家苑囿，时有大批民工陆续来此落户，村落进一步扩大，时称“千户屯”。永乐十二年（1414年）辟建南海子东、南、西、北四门，西红门为其西门，有兵丁驻守。

清康熙年间于此设置邮舍，多以“西红门”为名，逐渐取代了原来的“千户屯”，成为正式的村名。清末，南海子被废弃，西红门村逐渐发展成为京南重镇，为京都顺天府大兴县管辖，1958年划归北京市，今属西红门镇辖域。

西红门历史上以盛产萝卜而闻名京城。清光绪《顺天府志》记载：西红门“水萝卜，圆大如葖，皮肉皆绿，近尾则白。亦有皮红心白，或皮紫者，只可生食，极甘脆”。北京有句老话：“西红门萝卜叫城门。”明清时期北京城门黎明开，日暮闭。如果有宫里发的“腰牌”（特别通行证）便可以不按钟点随时进城，故称“叫城门”。据传有一年冬天，慈禧太后到南苑打猎，玩累了想吃梨，随行的太监想拿出早已准备好的梨，可是由于保暖食盒的盖子不知什么时候跑丢了，梨冻成了冰坨子。这时，西红门行宫的管事端上一盘切好的“心里美”萝卜，请“老佛爷”解渴，这才给太监救了急。慈禧太后见这萝卜翠绿的皮，紫红的心，透着一股鲜亮，只吃了两块儿，便觉得清脆爽口。她连声称好，并随即传旨，将西红门萝卜列为宫中贡品，按时令进奉。从此，只要西红门的菜农给宫里送萝卜，什么时候叫城门什么时候就开，于是便有了“西红门萝卜叫城门”之说。

\section{高米店北}
高米店北站位于大兴区北部，兴华大街与香园路交会处。

该站因地处前高米店村西北侧得名。《北京市大兴县地名志》载：“明初因村旁有一清河，故名清河店。清初因水患，清廷于此放赈赐米，改称‘高米店’。因聚落过大分设两村，有前后之分，该村称前高米店。清属顺天府大兴县黄村巡检司。1928年改属河北省大兴县黄村区，1949年属北平市24区，同年10月划属大兴县黄村区。1956年复属黄村镇至今。”

所谓“高米”，据传是高粱米的俗称。高粱米是高粱碾去皮层后的颗粒状成品粮。高粱又称红粮、蜀黍，古称蜀秫，是我国传统的五谷之一。

高米店具体何时发生水患及设置赈米发放机构，未见史料记载。民国时期，该村已成较大聚落，有几百户，村域呈矩形，主街道呈“井”字形。因地处永定河洪积—冲积平原，呈条状沙带地貌。2007年9月随着黄村新城北区建设，该村被整体拆迁，今属黄村镇辖域。

\section{高米店南}
高米店南站位于大兴区北部，兴华大街与金星西路交会处。

此站因地处黄村镇金星西路东端，初称“金星路站”，后因位于高米店村西南，且北部设有“高米店北站”而改称“高米店南站”。

另传，高米店曾是元代“下马飞放泊”（南苑）所建米粮仓。当时在园囿内种植水稻，一年一季，精耕细作，因“洼地种植，高地贮存”，俗称“高米”，是皇室内膳用米，所建储粮之仓称“高米殿”，后谐音为“高米店”。明代形成聚落仍称“高米店”，后分为两村，称“前高米店”“后高米店”，但此说无考。

民国时期已成较大聚落，有百余户，村域呈方形，街道网状分布。该村土壤系潮土沙土，经济以农业为主，农作物以小麦、玉米为主，还有蔬菜种植，兼种花生、大豆等油料作物。2014年设立高米店街道。

\section{枣园}
枣园站位于黄村镇西北部，枣园路与兴华大街交会处。

早年间该村处于永定河洪积—冲积平原，村落居于沙岗子之上，因村民多姓洪，故称“洪村岗子”，简称“洪村”。为固沙防风，村民们在沙岗子上种了许多枣树，明清时成枣林，其果实呈圆柱形，果皮红褐色，个大皮薄，核小肉多，质细甜脆，被称为洪村大枣、洪村脆枣，为大兴区特产，自明代开始已成皇室贡品，清嘉庆年间洪村大枣更是享誉京城。

据《黄村镇志》载：有一年秋天，慈禧太后去团河行宫，路过洪村歇歇脚，尝到了香美可口、又甜又酥的洪村大枣，很是高兴，便将此枣赐予外国使臣和朝中大臣，由此洪村大枣便有“贡枣”之称。此后每年到了中秋之后，宫里都要派人到大兴督办洪村大枣采摘事宜，以供慈禧太后品尝。

30多年前洪村一带还保留着枣树上百棵，1989年大兴县政府把这片枣林定名为“洪村枣园”，加以保护。此后在其附近建成住宅区，称“枣园小区”，东西走向的道路被称为“枣园路”，并建成“枣林公园”。

\section{清源路}
清源路站位于大兴区西北部，清源路与兴华大街交会处。

清源路为黄村卫星城内东西向主干线之一，西起兴业路，与清源西路相连，东至京开公路（今大广高速公路），与兴亦路（原清源东路）相接。其西段始建于1965年，旧名“黄鹅路”，即黄村至鹅房的县级公路。鹅房西邻永定河，历史上村域附近水多，村民多饲养鹅。清康熙皇帝巡视永定河堤时，常驻跸于此。相传康熙三十九年（1700年），一日河将决口，忽有一只大鹅自北飞来，引颈高歌，洪水遂退。康熙帝闻之，甚是惊喜，随即赐名“鹅房”。

“文革”时，“黄鹅路”曾改称“前进路”，1989年拓宽。东段形成于1979年，此后进行过改造。因大兴县第一自来水厂位于该路北侧，水厂建成后又称“水厂路”。1982年定名为“清源路”，寓意清泉水长流，饮水而思源。大兴县第一自来水厂是当时大兴县城供水面积最大的水厂，城区居民的主要生活用水多来自该厂，供水设备先进，输水管道密布。

与清源路相交的兴华大街，有黄村新城“第一街”之称，为城市主干道，北起北兴路，西南至义和庄，由原兴华北路、兴华中路和兴华南路改造而成，全线8.3公里，是黄村镇最为繁华的地段。

\section{黄村西大街}
黄村西大街站位于大兴区西北部，黄村西大街与兴华大街交会处。

黄村西大街为黄村卫星城内最早规划建设的商业街，建于1978年，初期曾名“向阳路”，长近600米，东西走向，东起兴丰大街，与黄村东大街相连，西至兴旺大街，并与兴华大街、兴业大街交会，因地处黄村镇西南侧而改称“黄村西大街”。

“黄村”这一地名始于元代，而成村于汉代，辽为析津县招贤乡东綦里。金元时这一带由于古浑河（今永定河）泛滥，风沙淤阻，形成大片沙丘荒地，出现聚落后遂称“荒村”，元代谐音为“黄村”，明代置黄村社。据传，元世祖忽必烈定都大都（北京）后，多次巡视京畿。一天他巡至大兴县的一个村落时，询问当地人：“此为何地？”当地人答称“荒村”，意谓荒凉之地。忽必烈觉得不雅，于是将“荒村”谐音为“皇村”，意为真龙王子所到之地，但当地人不明其意而谓之“黄村”，且沿用至今。

清康熙年间，黄村已成为京南商旅云集的重要集镇，直隶巡抚于成龙建议，在畿辅近地特设四路同知（知府的副职，正五品），其中南路同知府亦称南路厅。清代黄村已设官署，且建置完备，成为京城南部重镇。

自古以来，南方各省官员、商贾若陆路进京，多在此落脚，由此这里成为北京南部最大的驿站，素有“京师首驿”和“京南门户”之称。

\section{黄村火车站}
黄村火车站位于大兴区西北部，林校北路与兴华大街交会处。

该站因地处黄村火车站北侧得名。黄村火车站始建于清光绪十九年（1893年），两年后，即光绪二十一年（1895年），“京津铁路”筑成通车，投入使用，是北京历史上较为悠久的车站之一。黄村火车站的开通，为京南重镇黄村的繁荣起到了推进作用。

光绪二十五年（1899年）末，“义和团”起于山东，次年蔓延至直隶境内，包括今天的北京、天津地区。团民先是烧毁教堂、洋宅、施药房（由外国人开设的带有慈善性质的药房）、医院等，随后又将“洋人”修的铁路视为“洋玩意儿”进行破坏，扒掉铁路，烧毁电线，驱赶铁路维护工人，最后将破坏的重点放在火车站。北京周边的马家堡、长辛店、琉璃河、涿州、高碑店等火车站相继被毁，黄村火车站也未能幸免。“义和团”首先将黄村车站以及铁路沿线的几十根电线杆一起烧毁，使“京津铁路”陷入瘫痪状态。清末淮军名将聂士成奉命保护芦保、京津铁路，遭到“义和团”的迎击，伤亡数十人，随后黄村车站被焚毁，光绪二十八年（1902年）以后才得以修复。此后百余年，该站一直是京南铁路交通枢纽，曾有“京南第一站”之称。

如今的黄村火车站主体建于1992年，2011年进行了修葺，依然是京南铁路上的重要一站，京山线、京九线、西黄线等铁路线在此交会。

\section{义和庄}
义和庄站位于大兴区西北部，兴华大街与新源大街交会处。

义和庄明代属宛平县，清代属顺天府大兴县黄村巡检司，民国十七年（1928年）改属河北省大兴县黄村区，20世纪50年代划归北京市。据传，这一地名的由来与历史上的“义和团”有关。

“义和团”，又称“义和拳”，是清末发源于山东，针对西方在华人士所进行的大规模暴力运动，发端不久便受到各地的积极响应。“义和”二字出自《左传·襄公九年》：“利，义之和也。”即道义的总和。

清末时，受山东义和拳影响，有宛平县白家场、赵家场、王庄子等六个村子的村民自发组织起来设坛练拳，地点设在村中的一座庙宇内。为提高拳术，特从山东德州请来一位拳师为教习（教练）。当时这六个村子有百十来位二三十岁的青年都到庙宇里练拳。待拳术大大提高后，他们便参与了烧毁教堂、扒铁路等活动，同时将相邻的六个小村合并，改称“义和庄”，并铸造了一口大钟记录其事。

义和庄地处黄村镇西南五里许，东北距京山铁路约三里，北临佟家场，聚落略呈方形，有三条东西走向的街道，民国年间还只是河北省大兴县一个很小的聚落，村民居住较为分散，1949年以后随着村民的增加，村域不断扩大，20世纪90年代初为黄村镇行政村，有200余户，上千人。因地处永定河洪积—冲积平原，属条状沙带地貌，地下水位较浅，水量丰沛，历史上以农业生产为主，当时属大兴县黄村镇行政村。2009年该村整体拆迁，村民全部上楼。附近建起多处住宅区，西南临近念坛公园。

\section{生物医药基地}
生物医药基地站位于大兴区中西部，新源大街与永大路交会处南侧，京开高速公路西侧。

因北临中关村科技园区大兴生物医药基地得名，此地实称“韩园子”，故规划车站时初称“韩园子站”。

《北京市大兴县地名志》称：“明时村址原为纪百户庄韩家菜园子，清初有村民迁此，渐成村落，故名韩园子。清属顺天府宛平县。1928年改属河北省宛平县，1956年属大兴县黄村镇至今。”

“纪百户庄”在明代《宛署杂记》中已有载：“自（宛平）县出重城右安门，过大兴县界五十里，曰纪百户庄，分为二道。一道自纪百户庄东南十五里曰囤垡村，一道自纪百户正南七里曰天宫院村。”这里所说的“纪百户庄”，就是地处韩园子东北部的大庄村，如今已拆迁。《大兴史话》载：大庄明代成村，曾因纪姓以功封百户，得名纪百户庄。清朝末期，因永定河水患，该村受灾南移到今天的地址，因为人丁兴旺，村落较大，习惯称为大庄，1981年地名普查后正式定名为“大庄”。清光绪年间属顺天府大兴县黄村巡检司。

“百户”为古代官职，金初设置，为世袭军职。元代相沿，设百户为百夫之长，隶属于千户，而千户又隶属于万户，同为世袭军职。驻守各地者，设百户所，分隶于各千户所。据传明代此地为百户所军营，最高官员姓纪，称“纪百户”，其聚落称“纪百户庄”，后因此地有韩家菜园子而称“韩园子”。民国时期韩园子已成较大聚落，东北为大庄村，西南为张家庄，村域呈“L”形，主街道南北走向两条。因地处永定河洪积—冲积平原，土壤潮湿，农作物以小麦、玉米为主，兼种花生、蔬菜等。2002年以后因大兴生物医药基地的兴建，该村整体拆迁。

\section{天宫院}
天宫院站位于大兴区中西部，思邈路与新源大街交会处，京开高速路（大广高速路）西侧。

天宫院金代成村，原名“史家庄”。《宸垣识略》载：“天宫院在（宛平）县南六十里，旧名史家庄，因金章宗驾幸本村打围举膳，改名天宫院。”

关于“天宫院”的历史，史籍记载极少，据传其得名与金章宗有关。

金章宗完颜璟是金朝的第六位皇帝，女真名麻达葛，为金世宗完颜雍的嫡孙，在位20年，颇好巡游京畿。有一年春天，他带着十几个人到京南狩猎场打猎，不想马惊了，一路向南跑了四五公里，来到一个叫史家庄的村子。此时正是晌午，到了用膳时辰，一行人便走进路边的一家饭馆。章宗一边吃一边问随行的大臣，这村子叫什么名字。大臣也不知道，便到后厨问大师傅。大师傅说这村子叫史家庄，大臣刚要回禀皇上，只见章宗吃得正香，赶紧又回到后厨，告诉大师傅，千万别说这村子叫“史家庄”。大师傅问为什么，大臣说皇上正在用膳，提“史（屎）”字，太不吉利，属大不敬，是杀头的罪过。大师傅一听吓得连忙问道：“那叫什么呀？”大臣指着不远的一块空地问：“那块空地是干什么用的？”“是村民打场轧麦子的场院。”大臣想了想言道：“这村子就叫天宫院吧？天宫本是天帝、神仙居住之宫殿，皇帝贵为天子，用膳之地自是不凡之处。”大师傅应了一声。大臣随即向章宗回禀，说这村子叫“天宫院”，因早年间村中有座供奉天神的庙宇而得名。章宗信以为真，回宫后即兴御题了“天宫院”三个字。其字为“瘦金体”，颇有几分宋徽宗墨迹之遗韵。随后制成匾额，令人送到史家庄，从此“史家庄”就改称为“天宫院”。

明清时天宫院属宛平县，1958年划归北京市，曾属大兴区北臧乡，2009年7月设天宫院街道。

